\documentclass[11pt]{article}

\usepackage[T1]{fontenc}
\usepackage{amsmath}
\usepackage{amssymb}
\usepackage{hyperref}
\usepackage{tikz-cd}
\usepackage{tikz}
\usetikzlibrary{positioning}
\usepackage[titletoc,title]{appendix}

\tikzcdset{sep=huge}

\title{Chapter 05}
\author{Exercises}
\date{\today}

\begin{document}

\maketitle

\section{Exercise 1}

Exercise 1 Chapter 05: \textit{
  Show that a pullback of arrows
}
\[
\begin{tikzcd}
  A \times_{X} B \ar[r, "p_2"] \ar[d, "p_1"'] & B \ar[d, "g"] \\
  A \ar[r, "f"'] & X
\end{tikzcd}
\]
\textit{
in a category $\mathbf{C}$ is the same thing as their product in the slice
category $\mathbf{C}/X$.
}

\medskip

Objects of $\mathbf{C}/X$ are of shape $(A,f)$ with $f : A \to X$. Arrows
$h : (A,f) \to (B,g)$ are $h : A \to B$ with $g h = f$.

\textbf{Claim.} $(P,p_1,p_2)$ is a pullback of $f, g$ $\Leftrightarrow$
$\bigl((P,r),p_1,p_2\bigr)$ is a product of $(A,f), (B,g)$ in $\mathbf{C}/X$,
where $r = f p_1 = g p_2$. This is forced, since $p_1, p_2$ must be slice
arrows.

\textbf{Key fact ($\ast$).} If $f z_1 = s$, then every $v : Z \to P$ with
$p_1 v = z_1$ is a slice arrow $(Z,s) \to (P,r)$:
$r v = f p_1 v = f z_1 = s$.

\textbf{Solution.} We solve this by checking in either direction.

\begin{itemize}
  \item ($\Rightarrow$)
    Let $z_1 : (Z,s) \to (A,f)$, $z_2 : (Z,s) \to (B,g)$ be slice arrows, so
    $f z_1 = s = g z_2$. The pullback gives a unique $u : Z \to P$ with
    $p_i u = z_i$. By ($\ast$), $u$ is a slice arrow. Any other slice arrow
    with $p_i u' = z_i$ is an arrow of $\mathbf{C}$, so $u' = u$.
  \item ($\Leftarrow$)
    The square commutes: $f p_1 = r = g p_2$. Let $f z_1 = g z_2$ in
    $\mathbf{C}$ and put $s := f z_1$. Then $z_1, z_2$ are slice arrows out of
    $(Z,s)$, and the product gives a unique slice arrow $u$ with $p_i u = z_i$.
    Any $u' : Z \to P$ in $\mathbf{C}$ with $p_i u' = z_i$ is a slice arrow by
    ($\ast$), so $u' = u$. \hfill$\square$
\end{itemize}


\section{Exercise 2}

Exercise 2 Chapter 05: \textit{
  Let $\mathbf{C}$ be a category with pullbacks.
}

We use Exercise~1 and the \textbf{two-pullbacks lemma}: if in two adjacent
commuting squares the right one is a pullback, then the left one is a pullback
iff the outer rectangle is.

\subsection*{(a)}

\textit{Show that an arrow $m : M \rightarrow X$ in $\mathbf{C}$ is monic if
and only if the diagram below is a pullback.}
\[
  \begin{tikzcd}
    M \ar[r, "1_{M}"] \ar[d, "1_{M}"'] & M \ar[d, "m"] \\
    M \ar[r, "m"'] & X
  \end{tikzcd}
\]
\textit{Thus, as an object $\mathbf{C}/X$, $m$ is monic iff $m \times m \cong m$.}

\medskip

\textbf{Claim.} arrow $m: M \to X$ in $\mathbf{C}$ is monic
$\Leftrightarrow$ the diagram above is a pullback.

\textbf{Solution.} We solve this by checking in either direction:

\begin{itemize}
  \item($\Rightarrow$) Let $m z_1 = m z_2$. Since $m$ is monic, $z_1 = z_2$, so
    $u := z_1$ is the unique arrow with $1_M u = z_1$, $1_M u = z_2$.
  \item ($\Leftarrow$) Let $m z_1 = m z_2$. The pullback gives $u$ with $z_1 =
    1_M u = z_2$. So $m$ is monic.
\end{itemize}

\medskip
\textbf{As object of $\mathbf{C}/X$.} In $\mathbf{C}/X$ the arrow $m$ is the
object $(M,m)$, and $1_M$ is a slice arrow $(M,m) \to (M,m)$. By Exercise~1,
the square in (a) is a pullback iff $(M,m)$, with both projections $1_M$, is a
product of $(M,m)$ with itself in $\mathbf{C}/X$. This is what $m \times m
\cong m$ means: $m$ itself, with projections $1_M, 1_M$, is the product $m
\times m$. (Since products are unique up to iso, any other product of $m$ with
itself is then isomorphic to $m$, via $\langle 1_M, 1_M \rangle$.)

\pagebreak

\subsection*{(b)}

\textit{Show that the pullback along an arrow $f: Y \rightarrow X$ of a
pullback square over $X$,}
\[
  \begin{tikzcd}
    A \times_{X} B \ar[r] \ar[d] & B \ar[d] \\
    A \ar[r] & X
  \end{tikzcd}
\]
\textit{is again a pullback square over $Y$. (Hint: draw a cube and use the
two-pullbacks lemma.) Conclude that the pullback functor $f^{*}$ preserves
products.}

\medskip

\textbf{Solution.} Let $(P,p_1,p_2)$ be a pullback of $a : A \to X$, $b : B \to
X$, and $r := a p_1 = b p_2$. Pull back $a, b, r$ along $f$ to get $a' : A' \to
Y$, $b' : B' \to Y$, $r' : P' \to Y$ with top arrows $f_A, f_B, f_P$. The
induced $p_1' : P' \to A'$ satisfies $a' p_1' = r'$ and $f_A p_1' = p_1 f_P$
(it exists since $a(p_1 f_P) = r f_P = f r'$). $p_2'$ is constructed the same
way. So $a' p_1' = b' p_2'$.

\begin{enumerate}
  \item In
    \[
      \begin{tikzcd}[sep=large]
        P' \ar[r, "p_1'"] \ar[d, "f_P"'] & A' \ar[r, "a'"] \ar[d, "f_A"] & Y \ar[d, "f"] \\
        P \ar[r, "p_1"'] & A \ar[r, "a"'] & X
      \end{tikzcd}
    \]
    the right square and the outer rectangle are pullbacks, so the left square
    is a pullback, by the two-pullback lemma.
  \item In
    \[
      \begin{tikzcd}[sep=large]
        P' \ar[r, "p_2'"] \ar[d, "p_1'"'] & B' \ar[r, "f_B"] \ar[d, "b'"] & B \ar[d, "b"] \\
        A' \ar[r, "a'"'] & Y \ar[r, "f"'] & X
      \end{tikzcd}
    \]
    the right square is a pullback. The outer rectangle has the same arrows
    ($f_B p_2' = p_2 f_P$, $f a' = a f_A$) as the left square of Step~1,
    transposed, pasted with the original pullback, so it is a pullback. Meaning
    the left square is a pullback: $P' \cong A' \times_Y B'$.
\end{enumerate}

\textbf{$f^*$ preserves products.} $f^* : \mathbf{C}/X \to \mathbf{C}/Y$ sends
$(A,a)$ to $(A',a')$ and a slice arrow to the arrow induced by the pullback,
exactly as $p_1'$ above. So $f^*(P,r) = (P',r')$ and $f^*p_i = p_i'$. By
Exercise~1, $(P,p_1,p_2)$ is the product $(A,a) \times (B,b)$ in
$\mathbf{C}/X$ and  its image $(P',p_1',p_2')$ is a pullback over $Y$.
Therefore the product $f^*(A,a) \times f^*(B,b)$ is in $\mathbf{C}/Y$.

\subsection*{(c)}

\textit{Conclude from the foregoing that in a pullback square}
\[
\begin{tikzcd}
  M' \ar[r, "g"] \ar[d, "m'"'] & M \ar[d, "m"] \\
  A' \ar[r, "f"'] & A
\end{tikzcd}
\]
\textit{if $m$ is monic, then so is $m'$.}

\medskip
\textbf{Solution.} We use (a) to turn ``monic'' into a pullback square, (b) to
move that square along $f$, and (a) again to turn it back into ``monic''.

\begin{enumerate}
  \item \textbf{(by (a)).} Since $m$ is monic, the square
    \[
      \begin{tikzcd}
        M \ar[r, "1_M"] \ar[d, "1_M"'] & M \ar[d, "m"] \\
        M \ar[r, "m"'] & A
      \end{tikzcd}
    \]
    is a pullback.
  \item \textbf{(by (b)).} Pull this square back along $f : A' \to A$. Every
    corner $M$ lies over $A$ via $m$ (the top left one via $m 1_M = m$), and we
    may choose its pullback along $f$ to be the given square, so it becomes
    $M'$. The corner $A$ becomes $A'$. By (b) the resulting square
    \[
      \begin{tikzcd}
        M' \ar[r, "p_2'"] \ar[d, "p_1'"'] & M' \ar[d, "m'"] \\
        M' \ar[r, "m'"'] & A'
      \end{tikzcd}
    \]
    is a pullback.
  \item \textbf{(the legs are identities).} By construction, $p_1'$ is the
    unique arrow $x : M' \to M'$ with $m' x = m'$ and $g x = g$ (UMP of the
    given square). Since $1_{M'}$ satisfies both equations, $p_1' = 1_{M'}$.
    Likewise $p_2' = 1_{M'}$.
  \item \textbf{(by (a)).} So the square of Step 2 is exactly the square of (a)
    for $m'$, and it is a pullback. By (a), $m'$ is monic. \hfill$\square$
\end{enumerate}

\section{Exercise 3}

Exercise 8 Chapter 05: \textit{
  (Partial maps) For any category $\mathbf{C}$ with pullbacks, define the
  category $\textbf{Par}(\mathbf{C})$ of partial maps in $\mathbf{C}$ as
  follows: the objects are the same as those of $\mathbf{C}$, but an arrow $f:
  A \rightarrow B$ is a pair $(|f|,U_{f})$, where $U_{f} \rightarrowtail A$ is
  a subobject and $|f|: U_{f} \rightarrow B$ is a suitable equivalence class of
  arrows, as indicated in the diagram:
}
\[
\begin{tikzcd}
  U_{f} \ar[r, "\lvert f \rvert"] \ar[d, >->] & B \\
  A &
\end{tikzcd}
\]
\textit{
  Composition of $(\lvert f \rvert,U_{f}): A \rightarrow B$ and $(\lvert g
  \rvert,U_{g}): B \rightarrow C$ is given by taking a pullback and then
  composing to get $(\lvert g \circ f \rvert, \lvert f \rvert^{*}(U_{g}))$, as suggested by the
  following diagram:
}
\[
\begin{tikzcd}
  \lvert f \rvert^{*} (U_{g}) \ar[r] \ar[d, >->] & U_{g} \ar[r] \ar[d, >->] & C \\
  U_{f} \ar[r] \ar[d, >->] & B & \\
  A & &
\end{tikzcd}
\]
\textit{
  Verify that this really does define a category, and show that there is a
  functor,
  $$\mathbf{C} \rightarrow \textbf{Par}(C)$$
  which is the identity on objects.
}

\medskip
\textbf{Solution.} A representative of $f : A \to B$ is a pair $(m, |f|)$ with
a mono $m : U_f \rightarrowtail A$ and $|f| : U_f \to B$. Two pairs are
equivalent, $(m,|f|) \sim (m',|f|')$, iff there is an iso $i : U_f \to U_f'$
with $m' i = m$ and $|f|' i = |f|$. This is needed because pullbacks are only
unique up to iso: it makes composition well defined, and all equalities below
(identities, associativity, functor) hold as equivalence classes.

\begin{itemize}
  \item \textbf{Composition.} Pull back $U_g \rightarrowtail B$ along $|f|$ to
    get $V = |f|^*(U_g)$ with $n : V \to U_f$ and $v : V \to U_g$. Then $g
    \circ f = (|g| v,\ V \rightarrowtail U_f \rightarrowtail A)$.
    \begin{itemize}
      \item \emph{Domain is a subobject:} $n$ is a pullback of a mono, hence monic
        (Exercise~2(c)), and a composite of monos is monic.
      \item \emph{Well defined:} pullbacks are unique up to iso, and changing a
        representative by an iso changes the result only by an iso commuting with
        all arrows. So the class of $g \circ f$ does not depend on choices.
    \end{itemize}
  \item \textbf{Identities.} $1_A = (1_A, A)$. Pulling back along an identity,
    or pulling back an identity, gives the same object back. Meaning $f \circ 1_A
    = f$ and $1_B \circ f = (1_B |f|, U_f) = f$.
  \item \textbf{Associativity.} Let $h : C \to D$ be a third partial map, and
    let $T = |g|^*(U_h)$ be the domain of $h \circ g$. Pull back $T \to U_g$
    along $v$ to get $W$:
    \[
      \begin{tikzcd}[sep=large]
        W \ar[r] \ar[d] & T \ar[r] \ar[d] & U_h \ar[r, "|h|"] \ar[d, tail] & D \\
        V \ar[r, "v"] \ar[d, "n"'] & U_g \ar[r, "|g|"] \ar[d, tail] & C & \\
        U_f \ar[r, "|f|"'] \ar[d, tail] & B & & \\
        A & & &
      \end{tikzcd}
    \]
    All three squares are pullbacks. By the two-pullbacks lemma:
    \begin{itemize}
      \item pasting the top-left square with the one to its right shows that
        $W$ is the pullback of $U_h$ along $|g| v = |g \circ f|$, i.e.\ the domain
        of $h \circ (g \circ f)$.
      \item pasting it with the one below shows that $W$ is the pullback of
        $T \rightarrowtail U_g \rightarrowtail B$ along $|f|$, i.e.\ the domain of
        $(h \circ g) \circ f$.
    \end{itemize}
    Btivity,posites are $W \to V \to U_f \rightarrowtail A$ with
    $W \to T \to U_h \to D$, so they are equal.
\end{itemize}

\textbf{Functor.} Define $F : \mathbf{C} \to \mathbf{Par}(\mathbf{C})$
by $F(A) = A$ on objects and, for $f : A \to B$, $F(f) = (f, A)$: the partial
map defined on all of $A$, with domain $1_A : A \rightarrowtail A$.
\begin{itemize}
  \item \textbf{Identities:} $F(1_A) = (1_A, A)$ is the identity of
    $\mathbf{Par}(\mathbf{C})$.
  \item \textbf{Composition:} let $f : A \to B$ and $g : B \to C$. To compute
    $F(g) \circ F(f)$, pull back the domain $1_B$ of $F(g)$ along the map $f$
    of $F(f)$. The top-left square below is a pullback: if $z_1 : Z \to A$
    and $z_2 : Z \to B$ satisfy $f z_1 = z_2$, then $u := z_1$ is the unique
    arrow with $1_A u = z_1$ and $f u = z_2$. So $V = A$ with legs $1_A$ and
    $f$:
    \[
    \begin{tikzcd}
      A \ar[r, "f"] \ar[d, "1_A"'] & B \ar[r, "g"] \ar[d, "1_B"] & C \\
      A \ar[r, "f"'] \ar[d, "1_A"'] & B & \\
      A & &
    \end{tikzcd}
    \]
    Going right twice gives the map $g f$, going down twice gives the domain
    $1_A 1_A = 1_A$. So
    \[
      F(g) \circ F(f) = (g, B) \circ (f, A) = (g f, A) = F(g \circ f).
    \]
\end{itemize}
\hfill$\square$

\section{Exercise 4}

Exercise 12 Chapter 05: \textit{
  Consider the sequence of posets $[0] \rightarrow [1] \rightarrow [2] \rightarrow \dots$, where
  $$[n] = \{0 \leq \dots \leq n\},$$
  and the arrows $[n] \rightarrow [n + 1]$ are the evident inclusions.
  Determine the limit and colimit posets of this sequence.
}

\medskip
\noindent\textbf{Solution.}

Write $i_n : [n] \to [n+1]$, $k \mapsto k$. In $\mathbf{Pos}$ we compute the
(co)limit on underlying sets and then equip it with the order that makes the
legs monotone.

\medskip
\noindent\textbf{Colimit: $\omega = (\mathbb{N}, \leq)$.}
The union of all $[n]$ is $\mathbb{N}$, and $k \leq l$ holds in $[n]$ for any
$n \geq l$, so the order is the usual one. The legs are the inclusions
$c_n : [n] \to \mathbb{N}$; they satisfy $c_{n+1} i_n = c_n$.
Let $q_n : [n] \to Q$ be a cocone, i.e.\ $q_{n+1} i_n = q_n$. Define
\[
  u : \mathbb{N} \to Q, \qquad u(k) := q_n(k) \quad \text{for any } n \geq k .
\]
\begin{itemize}
  \item \emph{Well defined:} by compatibility, all $q_n$ with $n \geq k$ agree on $k$.
  \item \emph{Monotone:} for $k \leq l$ pick $n \geq l$; then
    $u(k) = q_n(k) \leq q_n(l) = u(l)$.
  \item \emph{Factors and unique:} $u c_n = q_n$ by definition, and every
    $k$ lies in some $[n]$, so $u$ is forced.
\end{itemize}
Note that $\omega$ is none of the $[n]$: it has no top element.

\medskip
\noindent\textbf{Limit: $[0] = 1$.}
An element of the limit is a family $(x_n)_n$ with $x_n \in [n]$ and
$i_n(x_n) = x_{n+1}$, i.e.\ $x_{n+1} = x_n$. Hence $x_n = x_0 \in [0] = \{0\}$
for all $n$, and the limit is the one-point poset with legs
$l_n : [0] \to [n]$, $0 \mapsto 0$.
Directly: a cone $a_n : L \to [n]$ satisfies $a_{n+1} = i_n a_n$, so it is
determined by $a_0 : L \to [0]$, which is the unique map to the terminal
object. So $[0]$ is a limit.

\medskip
\noindent\textbf{Why the asymmetry.} The index category $\omega$ has an
initial object $0$, so the limit is just the value $[0]$ there. It has no
terminal object, so the colimit is a new object, $\omega$.
\hfill$\square$

\end{document}
